\section{Discussion}
\label{sec:discussion}

\subsection{What Level 1.5 InEKF is, and what it is not}

Level~1.5 is a parallel Invariant Extended Kalman Filter on
$\mathrm{SO}(2)$, stacked on top of MapFormer's existing path-integration
scan. It adds three learnable components --- a constant prior covariance
$\Pi$, a per-token measurement noise $R_t$, and an inverse measurement
model $z_t$ --- connected by the scalar Hillis-Steele affine scan
\begin{equation}
  d_t = (1 - K_t) \, d_{t-1} + K_t \cdot \mathrm{wrap}(z_t - \hat\theta_{\mathrm{path},t}),
  \qquad
  \hat\theta_t = \hat\theta_{\mathrm{path},t} + d_t
\end{equation}
with Kalman gain $K_t = \Pi / (\Pi + R_t)$. The scan has the same
$\mathcal{O}(\log T)$ depth as MapFormer's underlying cumsum and the
same depth as Mamba's selective-scan primitive.

\paragraph{What it is.} A specialised 1-D state-space model with Kalman
parameterisation, Lie-group-correct wrapping, and input-dependent
measurement noise. The parameterisation makes $R_t$ interpretable as
per-token measurement uncertainty by construction --- its value has a
direct probabilistic meaning, not merely a learned scaling.

\paragraph{What it is not.} Level~1.5 is not a generic selective-scan
block that happens to work. Its structure is specifically chosen to
satisfy three requirements: (i) path-integration on $\mathrm{SO}(2)$
with the \citet{markovic2017wrapping} wrapped-innovation correspondence
to the Lie-Group EKF, (ii) state correction via an inverse measurement
model that respects the group structure, and (iii) compatibility with
MapFormer's existing geometric $\omega$ allocation. Removing any of
these via ablation degrades performance measurably
(Section~\ref{sec:results}).

\subsection{Relation to adjacent work}

\paragraph{MapFormer~\citep{rambaud2025mapformer}.} We reproduce
MapFormer faithfully (seven implementation details in
Section~\ref{sec:methods}) and inherit its input-dependent path
integration. Our contribution is orthogonal: it adds state correction
where MapFormer had none. On the paper's benchmark (aliased 2D torus
next-token prediction), both saturate at near-ceiling accuracy, so our
contribution lives in regimes the paper did not test --- noise,
landmarks, OOD length, calibration.

\paragraph{TEM / TEM-T~\citep{whittington2020tem,whittington2022tem-t}.}
TEM is the closest model family conceptually: it also learns
cognitive-map representations via action-conditioned recurrence. Our
per-visit-count curves (Section~\ref{sec:results:one-shot}) are the
TEM-canonical one-shot generalisation test, and Level~1.5 matches or
exceeds TEM's qualitative signature (sharp $k{=}1 \to k{=}2$ jump). TEM
trains on multiple environments simultaneously and compares
representations to neural data at length; we compare only at the
$\omega$-modular-spacing level, and candidly report negative results for
grid-cell and boundary-cell correspondence
(Section~\ref{sec:results:hippo}).

\paragraph{Mamba and selective SSMs~\citep{gu2024mamba}.} The connection
is deep. \citet{wang2025mamba-kalman} show Mamba's selective SSM is
isomorphic to a time-varying linear Kalman filter; Level~1.5's scan is
a 1-D Mamba block with three specific constraints (Kalman-ratio gain,
Lie-group wrap, decoupled path-integration + correction). The key
empirical finding is that this specialisation matters: our MambaLike
baseline (generic selective SSM at the same scale) fails to learn the
cognitive-map task (accuracy $\approx 0.57$ across configs), reproducing
the MapFormer paper's own Appendix A.5 Table~3 result. Generic SSMs can
in principle approximate this structure but cannot discover it from
gradient descent at our training scale; the explicit Lie-group prior
earns its keep.

\paragraph{Parallel Bayesian filters~\citep{sarkka2021parallel}.} Our
parallel InEKF is a specific instance of this broader program ---
specialised to $\mathrm{SO}(2)$ with a scalar affine scan rather than
the general matrix-associative primitive.

\paragraph{Contextual Positional Encodings.} CoPE~\citep{golovneva2024cope}
and TAPE~\citep{zhu2025tape} address a different problem: making
position encoding content-dependent at the token level. They do not
implement path integration on a compact group, and per the MapFormer
paper's Table~2 both fail on the 2D navigation task (CoPE 0.74 IID,
TAPE 0.23 IID vs MapFormer $\approx 1.00$). They are complementary
rather than competing; a CoPE-style content-based gating layered on top
of MapFormer could in principle provide landmark-detection gating for
the $R_t$ head (future work).

\subsection{Honest limitations}
\label{sec:discussion:limitations}

\paragraph{Single-environment training.} Each model trains on one fixed
\texttt{obs\_map} (paper-standard for MapFormer). Our OOD evaluations
use fresh obs\_maps and show ceiling-level zero-shot transfer to them,
but we do not do TEM-style multi-environment training where the model
sees thousands of distinct environments. The additional robustness such
training would provide is conjectural on our data; it is a natural
follow-up (Section~\ref{sec:future}).

\paragraph{Scale.} We follow MapFormer's 1-layer, 2-head, $d{=}128$
configuration. We have not scaled model or data meaningfully. A
comparison paper at 4 layers and 1M+ environments would change the story
in unknown ways.

\paragraph{Single Lie group.} Level~1.5 is developed for
$\mathrm{SO}(2)$, MapFormer's native group. The generalisation to
$\mathrm{SE}(2)$, $\mathrm{SE}(3)$, or higher dimensions involves the
same parallel-filter machinery but requires Adjoint-weighted
recurrences \citep{sola2018micro} rather than scalar affine scans. We
describe the path but have not implemented it.

\paragraph{Hippocampal correspondence.} As
Section~\ref{sec:results:hippo} reports, our current architecture does
not produce hexagonal grid cells (1 block per $\omega$ is structurally
insufficient for the three-waves-at-$60^\circ$ interference required),
and $R_t$ does not match the predicted Bayesian-optimal pattern.
Section~\ref{sec:future} proposes an architectural modification
(MapFormer-Grid) that could recover hexagonal organisation; testing
whether that actually works is left for follow-up.

\paragraph{Small Level15-EM variance.} One lm200 seed reached a
marginally worse final loss (1.40 vs $\sim 1.0$), inflating the lm200
std to $\pm 0.12$. A further-raised safe-init bias (5.0) partially
improves this seed but regresses others (the variance is intrinsic to
seed-specific landscape geometry, not a single-parameter issue). We
report this honestly rather than masking it.

\paragraph{CoPE baseline is incomplete.} Due to CoPE's high per-epoch
cost ($\sim 500$ min/run vs LSTM's 8 min/run on our setup), the lm200
CoPE row has 2 of 3 seeds. Our other baselines (MambaLike, LSTM, RoPE)
are all multi-seed complete.

\subsection{When to use Level 1.5 vs vanilla MapFormer}

Three rules of thumb emerge from our ablations and per-regime results:

\begin{enumerate}
  \item If the environment is noise-free and fully aliased, MapFormer
  is already at ceiling. Level~1.5's gain at $T{=}128$ clean is within
  noise of the baseline. Use Level~1.5 only if you need calibrated
  uncertainty (the NLL gap is still substantial).
  \item If actions are noisy or the environment has landmarks,
  Level~1.5 helps substantially (+11\,pp at $T{=}512$ OOD in both
  regimes). This is the primary use case.
  \item If you need calibrated uncertainty for downstream planning
  (e.g., active navigation), Level~1.5 is strictly preferable even on
  clean tasks.
\end{enumerate}

\subsection{Broader implication: Mamba-as-Kalman on a Lie group}

A broader reading of our work: Level~1.5 is a specialised Mamba-style
selective-scan on $\mathrm{SO}(2)$. The \citet{wang2025mamba-kalman}
view (Mamba $\equiv$ linear Kalman filter) tells us that selective SSMs
\emph{can} implement Kalman filters. Level~1.5 shows that they
\emph{should}, at least on tasks with known geometric structure,
because the specialisation (Kalman-ratio gain, Lie-group wrapping,
decoupled path-integration scan) yields a large improvement in what the
model can discover from gradient descent at limited training scale.

This suggests a broader research direction: \emph{constrained Mamba
architectures with explicit group-equivariance}. $\mathrm{SO}(2)$ is
the simplest case; $\mathrm{SE}(2)$ and $\mathrm{SE}(3)$ are natural
generalisations for 2D- and 3D-navigation tasks; other compact Lie
groups cover other periodic structures. In each case, the specialised
Mamba block would look like Level~1.5 but with an Adjoint-weighted scan
and higher-dimensional measurement model. We expect the gradient-descent
discoverability advantage to persist across all of them: \emph{generic
selective SSMs cannot discover geometric structure at small scale;
explicit group-aware selective SSMs can.}

\subsection{Conclusion}

We started from a reproduction of MapFormer, added a parallel Invariant
EKF machinery preserving its $\mathcal{O}(\log T)$ parallelism, and
characterised the combined architecture across three untested regimes.
The main concrete contribution --- Level~1.5 --- provides bounded-error
length generalisation, robustness to action noise, strong landmark
utilisation, and calibrated uncertainty, at the cost of a single
learnable $\Pi$, one per-token MLP for $R_t$, one inverse measurement
MLP, and a scalar affine scan. It is mechanistically simple,
architecturally minimal, and empirically effective.

We are not claiming to solve a task MapFormer could not. We are
claiming that classical Lie-group filtering, parallelised, is a natural
extension of MapFormer-family architectures into precisely the regimes
that any deployed navigation system must handle. The transfer direction
--- classical robotics $\to$ deep learning, via explicit group-aware
primitives --- remains underexplored, and this paper is one step along
it.
